% Claim statement file for row claim_3_27 -- "LabelRoman".
% Auto-generated stub by scaffold/solve_chapter.py at row start. The
% manager agent (or a worker it dispatches) fills the body below with the
% claim's statement(s) from the lecture notes -- statement only, no proof
% (proofs live in the sibling `claim_3_27_proof_*.tex` file).
%
% This file is a sibling of `main.tex` inside `<subsection>/tex/`. The
% `subfiles` package lets it render both standalone and as part of
% `main.tex`. Note: `main.tex` does NOT \subfile this statement file -- the
% sibling `_proof_` file restates the statement above the proof, so
% including both in `main.tex` would be redundant. This file exists for
% standalone reading / grep convenience.

\documentclass[main]{subfiles}

\begin{document}
% ref: claim_3_27 (statement)
% title: LabelRoman
\def\rowref{claim\_3\_27}
\phantomsection\label{claim_3_27}

\begin{Lem}[LabelRoman]\label{lem:replace_walk}
    Let $G = (J, V, E, L)$ be a CDMG (def \ref{def-cdmg}); throughout we use the notation of def \ref{not-cdmg}, the walk and directed-walk definitions of def \ref{def:walks}, items~i.\ and~ii., the family-relationships definition of def \ref{def_3_5} (specifically the ancestor set $\Anc^G(v) \subseteq J \cup V$, the descendant set $\Desc^G(v) \subseteq J \cup V$, and the strongly-connected-component set
    \[
        \Sc^G(v) \;:=\; \Anc^G(v) \cap \Desc^G(v) \;\subseteq\; J \cup V
    \]
    for each $v \in J \cup V$, with $v \in \Sc^G(v)$ by reflexivity of $\Anc^G$ and $\Desc^G$), the ancestor set of a subset $\Anc^G(C) := \bigcup_{c \in C} \Anc^G(c)$ from def \ref{def_3_5}, item~iv.\ (and as used in def \ref{def:sigma_blocking}), the collider/non-collider machinery of def \ref{def:collider_noncollider}, the blockable/unblockable refinement of def \ref{def:unblockable_noncollider}, and the $\sigma$-open/$\sigma$-blocked classification of walks of def \ref{def:sigma_blocking}.

    Let $C \subseteq J \cup V$ be a subset of nodes; $C$ is \emph{not} assumed disjoint from $J$, nor non-empty. Let $n \ge 0$ be an integer and let
    \[
        \pi \;=\; (v_0, a_0, v_1, a_1, \dots, v_{n-1}, a_{n-1}, v_n)
    \]
    be a walk in $G$ in the sense of def \ref{def:walks}, item~i., with $v_0, \dots, v_n \in J \cup V$ and $a_0, \dots, a_{n-1} \in E \cup L$ satisfying the walk constraint at every $k \in \{0, 1, \dots, n - 1\}$:
    \[
        \bigl(a_k = (v_k, v_{k+1}) \;\land\; a_k \in E \cup L\bigr)
        \;\;\lor\;\;
        \bigl(a_k = (v_{k+1}, v_k) \;\land\; a_k \in E\bigr).
    \]
    Assume that $\pi$ is $C$-$\sigma$-open in the sense of def \ref{def:sigma_blocking}.

    \emph{Hypothesis on positions $(i, j)$.}
    Suppose $i, j \in \{0, 1, \dots, n\}$ are positions on $\pi$ with $i < j$ and
    \[
        v_i \in \Sc^G(v_j) \quad\bigl(\text{equivalently, } v_i \in \Anc^G(v_j) \;\land\; v_i \in \Desc^G(v_j)\bigr).
    \]
    The strict inequality $i < j$ orders the \emph{positions} on $\pi$ but does \emph{not} constrain the underlying \emph{vertices}: per def \ref{def:walks}, item~i., ``the same node may appear multiple times in a walk'', so the special case $v_i = v_j$ is admissible (and is exactly the case in which $v_i \in \Sc^G(v_j) = \Sc^G(v_i)$ is automatic by the reflexivity $v_i \in \Sc^G(v_i)$ of def \ref{def_3_5}).

    \emph{Shortest directed walk between two nodes.}
    For nodes $u, w \in J \cup V$ with $u \in \Anc^G(w)$, a \emph{shortest directed walk from $u$ to $w$ in $G$} is a directed walk from $u$ to $w$ in $G$ (def \ref{def:walks}, item~ii.) of minimum length among all directed walks from $u$ to $w$ in $G$. Spelled out: it is a finite sequence
    \[
        (u_0, b_0, u_1, b_1, \dots, b_{m-1}, u_m)
    \]
    for some integer $m \ge 0$, with $u_0, \dots, u_m \in J \cup V$, $u_0 = u$, $u_m = w$, and $b_l = (u_l, u_{l+1}) \in E$ for every $l \in \{0, 1, \dots, m - 1\}$ (so every walk-edge is a directed edge whose tail is the earlier vertex and whose head is the later vertex on the sequence), such that no directed walk from $u$ to $w$ in $G$ has length strictly less than $m$.

    The hypothesis $u \in \Anc^G(w)$ ensures the set of directed walks from $u$ to $w$ in $G$ is non-empty (def \ref{def_3_5}, item~iv.), so the well-ordering of $\mathbb{N}$ yields a witness of minimum length. The degenerate case $u = w$ is admitted at $m = 0$: the single-vertex tuple $(u_0) = (u)$ is a directed walk from $u$ to $u$ in $G$ in the sense of def \ref{def:walks}, item~ii.\ at $n = 0$ (no walk-edges to constrain), and is therefore the unique shortest directed walk from $u$ to $u$ in $G$ (every other directed walk from $u$ to $u$ in $G$ has length $\ge 1$).

    \emph{Replacement subwalk $\sigma_{ij}$ and modified walk $\pi'$.}
    Define a finite alternating sequence
    \[
        \sigma_{ij} \;=\; (w_0, b_0, w_1, b_1, \dots, b_{m-1}, w_m)
    \]
    for some integer $m \ge 0$ (with $m$ and the data $w_0, \dots, w_m \in J \cup V$, $b_0, \dots, b_{m-1} \in E \cup L$ determined by the case below), satisfying $w_0 = v_i$ and $w_m = v_j$, by exactly one of the following two mutually-exclusive jointly-exhaustive cases on the boundary configuration of $\pi$ at position $j$:

    \begin{enumerate}[label=(\roman*)]
        \item \textbf{Case (i):} $j = n \;\;\lor\;\; a_j = (v_j, v_{j+1}) \in E$.

            (The first disjunct is $j = n$; the second disjunct requires $j \le n - 1$ and asserts that the walk-step $a_j$ is the \emph{directed} edge from $v_j$ to $v_{j+1}$ — equivalently, $a_j$ is an edge in $E$, $a_j$ is \emph{not} an edge into $v_j$ in the sense of def \ref{def:collider_noncollider}, and $a_j$ \emph{is} an outgoing walk-edge of $v_j$ at position $j$ on $\pi$ in the sense of def \ref{def:unblockable_noncollider}.)

            Let $\sigma_{ij}$ be a shortest directed walk from $v_i$ to $v_j$ in $G$:
            \[
                (w_0 = v_i, \; b_0, \; w_1, \; b_1, \; \dots, \; b_{m-1}, \; w_m = v_j),
                \qquad
                b_l = (w_l, w_{l+1}) \in E \;\;\text{for every } l \in \{0, 1, \dots, m - 1\},
            \]
            of minimum length over all directed walks from $v_i$ to $v_j$ in $G$. Existence of such a witness uses $v_i \in \Sc^G(v_j) \subseteq \Anc^G(v_j)$ (def \ref{def_3_5}). When $v_i = v_j$ the minimum length is $m = 0$ and $\sigma_{ij} = (v_i)$ (single vertex, no walk-edges).

        \item \textbf{Case (ii):} $\lnot\,$(case~(i)), i.e.\ $j \le n - 1 \;\;\land\;\; \lnot\bigl(a_j = (v_j, v_{j+1}) \in E\bigr)$.

            (Spelled out via the walk constraint at index $j$: $j \le n - 1$ and either $a_j = (v_j, v_{j+1}) \in L$, or $a_j = (v_{j+1}, v_j) \in E$. Equivalently: $a_j$ is an edge \emph{into} $v_j$ in the sense of def \ref{def:collider_noncollider}'s ``edge into $v_k$'' predicate at $k = j$, and $a_j$ is \emph{not} an outgoing walk-edge of $v_j$ at position $j$ on $\pi$ in the sense of def \ref{def:unblockable_noncollider}.)

            Let $\sigma_{ij}$ be the reversal of a shortest directed walk from $v_j$ to $v_i$ in $G$. Concretely, fix integers $m \ge 0$ and vertices $u_0, \dots, u_m \in J \cup V$ such that the sequence
            \[
                (u_0 = v_j, \; c_0, \; u_1, \; c_1, \; \dots, \; c_{m-1}, \; u_m = v_i),
                \qquad
                c_l = (u_l, u_{l+1}) \in E \;\;\text{for every } l \in \{0, 1, \dots, m - 1\},
            \]
            is a directed walk from $v_j$ to $v_i$ in $G$ of minimum length over all directed walks from $v_j$ to $v_i$ in $G$. Existence of such a witness uses $v_i \in \Sc^G(v_j) \subseteq \Desc^G(v_j)$ (def \ref{def_3_5}, applied in the form $v_j \in \Anc^G(v_i)$).

            Set $w_l := u_{m - l}$ for $l \in \{0, 1, \dots, m\}$ and $b_l := c_{m - 1 - l}$ for $l \in \{0, 1, \dots, m - 1\}$, so that $\sigma_{ij} = (w_0 = v_i, b_0, w_1, \dots, b_{m-1}, w_m = v_j)$ and every walk-edge satisfies
            \[
                b_l \;=\; c_{m - 1 - l} \;=\; (u_{m - 1 - l}, u_{m - l}) \;=\; (w_{l+1}, w_l) \;\in\; E
                \qquad\text{for every } l \in \{0, 1, \dots, m - 1\}.
            \]
            (So each $b_l$ is a directed edge in $E$ whose head is $w_l$ and whose tail is $w_{l+1}$ — the second disjunct of the walk constraint of def \ref{def:walks}, item~i., at position $l$ on $\sigma_{ij}$. The sequence $\sigma_{ij}$ is \emph{not} itself a directed walk from $v_i$ to $v_j$ in the sense of def \ref{def:walks}, item~ii.; it is the walk-reversal of the directed walk $(u_0, c_0, \dots, c_{m-1}, u_m)$ from $v_j$ to $v_i$.) When $v_i = v_j$ the minimum length is $m = 0$ and $\sigma_{ij} = (v_i)$ (single vertex, no walk-edges).
    \end{enumerate}

    The two cases are mutually exclusive and jointly exhaustive over the conditions on $(j, a_j)$ given $j \le n$: case~(i) covers ``$j = n$, or $j \le n - 1$ and $a_j = (v_j, v_{j+1}) \in E$''; case~(ii) covers ``$j \le n - 1$ and the negation of $a_j = (v_j, v_{j+1}) \in E$''. The case-(i)/case-(ii) classification depends on the boundary configuration of $\pi$ at position $j$ only, not on the boundary configuration at position $i$.

    Define the \emph{modified walk} as the concatenation of the prefix $(v_0, a_0, \dots, a_{i-1}, v_i)$ of $\pi$, the replacement subwalk $\sigma_{ij}$ above, and the suffix $(v_j, a_j, v_{j+1}, \dots, a_{n-1}, v_n)$ of $\pi$:
    \[
        \pi' \;=\; \bigl(v_0, \, a_0, \, v_1, \, \dots, \, a_{i-1}, \, v_i = w_0, \, b_0, \, w_1, \, b_1, \, \dots, \, b_{m-1}, \, w_m = v_j, \, a_j, \, v_{j+1}, \, \dots, \, a_{n-1}, \, v_n\bigr).
    \]
    Then $\pi'$ is a walk in $G$ in the sense of def \ref{def:walks}, item~i., of length $i + m + (n - j)$: the prefix segment $(v_0, a_0, \dots, a_{i-1}, v_i)$ inherits the walk constraint at every $k \in \{0, 1, \dots, i - 1\}$ from $\pi$; the suffix segment $(v_j, a_j, v_{j+1}, \dots, a_{n-1}, v_n)$ inherits the walk constraint at every $k \in \{j, j + 1, \dots, n - 1\}$ from $\pi$; the replacement segment $(w_0, b_0, \dots, b_{m-1}, w_m)$ satisfies the walk constraint at every $l \in \{0, 1, \dots, m - 1\}$ — by the first disjunct $b_l = (w_l, w_{l+1}) \in E$ in case~(i), and by the second disjunct $b_l = (w_{l+1}, w_l) \in E$ in case~(ii).

    \emph{Conclusions.}
    \begin{enumerate}
        \item \textbf{The replacement subwalk lies entirely in $\Sc^G(v_j)$:}
            \[
                \forall\, l \in \{0, 1, \dots, m\}, \quad w_l \in \Sc^G(v_j).
            \]
            (Both endpoints — $w_0 = v_i \in \Sc^G(v_j)$ by the hypothesis on $(i, j)$, and $w_m = v_j \in \Sc^G(v_j)$ by reflexivity — and every interior vertex $w_l$ for $l \in \{1, 2, \dots, m - 1\}$, the last being vacuous when $m \le 1$ and in particular when $m = 0$.)
        \item \textbf{The modified walk $\pi'$ is $C$-$\sigma$-open} in the sense of def \ref{def:sigma_blocking}.
    \end{enumerate}

    \emph{Addition to the LN: length-zero replacement when $v_i = v_j$.}
    The hypothesis $i < j$ pins down the positions on $\pi$ but not the underlying vertices, so the case $v_i = v_j$ is admissible whenever $\pi$ revisits a node. In that case $v_i \in \Sc^G(v_j) = \Sc^G(v_i)$ holds by the reflexivity $v_i \in \Sc^G(v_i)$ of def \ref{def_3_5} alone, and the prescribed ``shortest directed walk from $v_i$ to $v_j$ in $G$'' (case~(i)) or ``shortest directed walk from $v_j$ to $v_i$ in $G$'' (case~(ii)) is the length-$0$ trivial directed walk: the single-vertex tuple $(v_i) = (v_j)$ with empty edge list, per def \ref{def:walks}, item~ii.\ at $n = 0$. Any positive-length directed walk from $v_i$ to $v_i$ in $G$ has length $\ge 1$ and is therefore not shorter than the length-$0$ trivial directed walk, so the trivial directed walk is unconditionally the unique minimum (and in particular the existence of a directed self-loop $(v_i, v_i) \in E$, or of longer directed cycles through $v_i$, does not affect the minimum-length witness selected here). Under this collapse the replacement subwalk is $\sigma_{ij} = (v_i)$ (one vertex, no walk-edges), every vertex of $\sigma_{ij}$ lies in $\Sc^G(v_j)$ trivially, and the modified walk $\pi'$ has length $i + 0 + (n - j) = n - (j - i) < n$ — strictly shorter than $\pi$ (the $j - i$ walk-steps of $\pi$ between positions $i$ and $j$ are dropped). The Lean formalisation admits length-$0$ directed walks uniformly (the single-vertex tuple satisfies the directed-walk constraint of def \ref{def:walks}, item~ii.\ vacuously), so this degenerate branch is included in the statement above without a separate carve-out: the existential ``$\exists \sigma_{ij}, \pi'$'' in the conclusion is witnessed by the length-$0$ subwalk and the suitably-shortened modified walk in this case, and conclusion~(1) (every vertex of $\sigma_{ij}$ in $\Sc^G(v_j)$) and conclusion~(2) ($\pi'$ is $C$-$\sigma$-open) are both met.
\end{Lem}

\end{document}
